\documentclass[11pt]{article} 
\usepackage{nopageno}
\setcounter{secnumdepth}{1} % Use 1 for sections, 2 for subsections, etc.
\usepackage[margin=1in]{geometry}
\begin{document}

% Remove page numbers---Research.gov will paginate the document instead
%
\pagestyle{empty} 

\noindent \textbf{Facilities, Equipment, and Other Resources:}

\begin{itemize}

\item \textbf{The Population Research Institute (PRI) at Penn State:} 
The Population Research Institute (PRI) at Penn State occupies over 6800 square feet of office space, meeting facilities, and computer labs available for use by faculty and postdoctoral associates, staff, and research assistants. The primary physical facilities of PRI are centrally located on the Penn State University Park campus in similarly close proximity to most participating academic departments and colleges. The physical facilities include open space that houses equipment for copying, scanning, and large-quantity printing and a mailroom with adjacent reception area for receiving visitors. All staff members have up to date desktop computers and printers in their offices with sufficient space for meetings with investigators. Shared meeting space is available for larger meetings. Beyond these resources available to all of PRI, each core is well equipped to manage and support associates’ research activities according to each core’s own specific aims. Staff in each core are able to provide services to associates working remotely via shared network file storage, box file sharing and virtual conferencing.  

 

PRI’s Administrative core manages facilities including the allocation of office space, furnishings, building maintenance, equipment, keys, card-swipe access into computer labs and restricted data enclaves, and scheduling for conference rooms. All staff in the Administrative core have offices in or adjacent to the main reception area (three offices). The Director has an office in this same suite and the Associate Director’s office is next door to the suite. This proximity ensures that the Core can efficiently provide pre- and post- award support, financial services, compliance with regulations of the University and funding organizations, and PRI sponsored events and communications. The Core also oversees office assignments to PRI research associates, postdoctoral researchers, and visiting scholars.  

 

PRI has two dedicated conference rooms, one with capacity for twenty people and one with capacity for eight. These rooms include state of the art technology to facilitate communication including shared screen presentations with colleagues in person and in remote locations. The larger conference room (605) is equipped with one 65 inch TV/Monitor, conference camera, sound bar (speaker), multi-directional room mic, a long HDMI cable connected to the TV for alternate input (instead of the podium PC) and USB-C, DisplayPort, Mini DisplayPort, and VGA video adapters for the HDMI cable. The smaller room located one floor above (707) is also equipped for in-person and remote presentations. These facilities promote communication with research collaborators across Penn State and around the world.  

 

PRI computer services are managed by the SSRI Research IT group. The PRI computer lab, maintained by SSRI IT staff, contains 12 workstations. The lab was recently upgraded with Dell Optiplex 5050’s running Windows 10. The lab maintains support for the latest versions of Stata, SPSS, SAS, R, other data analysis and statistical software, as well as Microsoft Office and the Adobe suite of products. Another application server supports ArcGIS accessible from the computer lab or secure VPN. Researchers also use the lab to access the PRI servers for approved restricted data sets that are permitted to be accessed this way. In addition, PRI associates and researchers have access to an application server dedicated to research computing. This server is a Dell PowerEdge R730 running Windows Server 2019 with 512G of memory. It contains the latest versions of the same software available in the computer lab and single licenses of more expensive software such as Atlas.ti and Mplus. In addition, the server has the multi-processor version of Stata, which can run large jobs more efficiently. This server is accessible anywhere in the world via the University VPN.  

 

PRI has restricted servers to secure access to data. For example, one server is dedicated to Add Health via a virtual machine running Windows Server 2019 and located in the University Park data center in a facility proximate to the main campus at University Park, PRI’s physical home. PRI hosts eight additional restricted data enclaves for datasets requiring stand-alone non-networked computers. The largest enclave is dedicated to data products from NCES. Supervision for enclave security is maintained by PRI CSA compliance staff who work closely with SSRI Research IT to maintain and audit security and software. Other restricted datasets reside on non-networked computers in separate offices assigned to PRI investigators as needed. Computers for these datasets are encrypted with BitLocker and require a password to log in. 

 

Additional equipment (i.e. laptops, projectors, and tablets) for short term use and support for remote access, online webinar and meeting hosting are all available in PRI from SSRI Research IT, services that were tremendously useful during the most recent period of campus restrictions during the coronavirus pandemic. Although in-person, face to face collaboration is often preferred, PRI has the capability to continue its activities remotely and keep associates engaged with SSRI IT support for online meetings, methodology workshops, webinars, and the NIH writing group when necessary.  

\item \textbf{Penn State Federal Statistical Research Data Center (FSRDC):} 
This center is one of 30 Federal Statistical Research Data Centers in the United States. The Penn State FSRDC is located in secured rooms in the main library on the Penn State University Park Campus – a central location that is within easy reach of faculty and researchers and the central offices of PRI. The FSRDC has since expanded with partnerships across Pennsylvania including satellites in Philadelphia and Pittsburgh. The FSRDC provides secure access to restricted economic, demographic, and health data collected by U.S. Census Bureau, the National Center for Health Statistics, the Agency for Healthcare Research and Quality, and the Bureau of Labor Statistics. It is a vital resource for Penn State faculty and researchers in the fields of economics, business, demography, statistics, sociology, and health services. Access to the FSRDC requires a research project approved by the Census Bureau or agency that collects the data. Researchers with approved projects receive Special Sworn Status from the Census Bureau and are able to work within the data center. PRI associates have access to statistical programming support from CSA and pre-award support from Administrative core staff for assistance preparing applications for approval to access the FSRDC.  

\item \textbf{The Penn State Administrative Data Accelerator (ADA):} 
The Penn State Administrative Data Accelerator’s PA Integrative Data System aims to supports partnerships between researchers and state government through the use of administrative records to expedite research on critical social issues. To accomplish this, the Accelerator links datasets in a highly secure environment—leveraging analytic and privacy best practice. This infrastructure facilitates engagement with policymakers to translate research into evidence-based policy and budgetary insights. In addition to providing support for the Penn State research community’s access and use of administrative data, the accelerator maintains a secure NIST 800-171 ATO approved enclave for holding sensitive records.  

\item \textbf{Heritage 3D lab at PSU:}
The Penn State Heritage 3D Lab, led by PI Rebecca Napolitano, is a state-of-the-art research facility dedicated to advancing the understanding and preservation of heritage structures through innovative engineering and computational methods. The lab brings together a diverse team of faculty, graduate and undergraduate students to pioneer research on structural analysis, heritage preservation, and disaster resilience. The H3D Lab houses a wide range of field-deployable equipment supporting both research and teaching, including a Trimble X7 terrestrial laser scanner, Leica laser rangefinder, panoramic photosphere camera, various environmental sensors, and thermal imaging cameras. These tools enable on-site assessments of heritage structures and their environments.

The lab also support a broad range of research related to built environment visualization, automation, robotics, and the application of advanced sensing technologies for health and safety monitoring on heritage sites and in infrastructure management. 
Key equipment includes two advanced mobile UGV robots equipped with multiple sensors and a robotic manipulator arm, as well as advanced wearable and immersive technologies such as AR and VR headsets for model visualization and digital twin interactions.
The lab also utilizes various wearable sensors for real-time monitoring of workers' mental and physical health during field operations.

The lab also features an advanced workspace designed for integrated, team-based workshops. The ICon Lab features an immersive curved screen with ultra-high-definition projection capabilities, supporting 3D visualization and walkthroughs of architectural models. This space facilitates collaborative design, review, and training sessions for up to 30 people, pushing interdisciplinary teams together.

\item \textbf{Busse Lab at PSU:}
The Busse lab, led by PI Busse, currently has four Dell Precision 3660 desktop computers dedicated to computational work. These computers are equipped with SimaPro for life cycle assessment (LCA), Ecoinvent for LCA results interpretation, Solidworks, EPANET, python, HEC-RAS 2D, and Matlab. 

The main lab space is equipped for water treatment and water system testing applications. This space is approved to conduct BSL2 level experiments. 

Dr. Busse is one of three faculty that shares a closed loop water channel built by Engineering Laboratory Design, Inc.

\item \textbf{Center for Socially Responsible Artificial Intelligence (CSRAI) at Penn State:} 
PI Napolitano, an affiliate member of CSRAI, will leverage the Center's expertise in promoting, practicing, and studying socially responsible ways of building, deploying, and using AI technology, with an emphasis on understanding the social and ethical implications of these activities. CSRAI's holistic approach, which encourages all AI research and development activities to consider social and ethical implications as well as intended and possible unintended consequences, will inform the project's efforts to develop a framework for responsible disaster management and cultural heritage preservation. The Center's multi-disciplinary research in AI and its focus on involving all stakeholders, including experts in science, technology, humanities, business, policy, and law, will support the project's goal. 

\item \textbf{Survey Research Center (SRC)}: The Survey Research Center (SRC) at Penn State offers fee-based services to design and implement a variety of data collection methods. SRC provides infrastructure for Computer Assisted Telephone Interviewing (CATI), support for traditional mail surveys (i.e., design, print, and scan self-administered surveys, and services for envelope stuffing, mailing, and survey receiving services). SRC also supports cost-effective internet surveys written with Qualtrics or RedCap software. SRC staff also are available to support focus groups and in-person fieldwork. Finally, SRC supports the Dynamic Real-Time Ecological Ambulatory Methodologies (DREAM) program for ecological momentary assessment, experience sampling, daily diary techniques, ambulatory monitoring of physiological function, physical activity, and/or movement monitoring, and the assessment of ambient environmental parameters.  

\item \textbf{Teleconference and Web Meeting Facility:} Penn State provides teleconference capability with a video option that is free to the research team when any team member is traveling. We will use the teleconference capability and web meeting facility for webinars through which we will disseminate our research outcomes to other researchers, governments, businesses, and non-profit organizations.  

\item \textbf{Institute for Computational and Data Sciences-Advanced Cyber Infrastructure:}
Over the course of the research project, research data will be hosted by the Pennsylvania State University’s Institute for Computational and Data Sciences (ICDS) through its Roar supercomputer. Roar provides both active storage (for data that is being worked on, requiring frequent access) and near-line storage (for back-up purposes and data that needs only infrequent access). Active storage is achieved through DDN 12KX40 and GS7K flash storage array systems, while near-line storage utilizes Oracle’s FS1 flash storage appliance and a SL8500 Tape Library. Active storage is backed up to the SL8500 daily. 
We will also use Roar to archive the research data for at least three years after the end of the award or after the public release of the data, whichever comes later. ICDS provides long-term archival services using the Oracle SL8500. Multiple copies of archived data are created through Oracle Hierarchical Storage Manager (Oracle HSM) to safeguard against data corruption. Oracle HSM generates and maintains metadata on archived files so that the data can be readily accessed. This technology affords us easy retrieval of data, even years after it has been written. 
Data Security: ICDS implements various security measures to ensure that data stored on the Roar system remains safe. Roar requires a strong password and two-factor authentication for access, and all access can be audited by ICDS staff. To mitigate the potential for malicious software and security attacks, Roar employs automated weekly scans for identifying and patching software vulnerabilities. Roar provides the capability to encrypt data in-flight (when moving between points) and at rest (while written in storage). Roar login/endpoint nodes are protected by software-based firewalls that only permit Secure Shell (SSH) traffic. By default, ICDS enforces “Least Privilege” access concepts across the system, providing users with only the minimum set of permissions and accesses required to complete their function. Roar storage is physically protected in Penn State’s Tower Road Data Center (TRDC). Physical access to the systems is limited to systems administration personnel with exceptions controlled by the TRDC’s secure operations center. TRDC requires swipe-card access and is monitored at all times. Data stored on Roar’s active storage systems is backed up to tape storage for a period of 90 days. Backup data is automatically purged from tape once the 90 days has been exceeded. 

\item \textbf{PSU Career Services:}
The Graduate School and Graduate Career Services offer an array of programs, resources, and services that can support professional and career development. Services offered by Lesley Jackson, graduate student career and professional development, will include career counseling, mock interviewing, group counseling and more. Graduate student programs are offered throughout the academic year and in different formats.  While many different programs are presented, frequent topics offered include Career Exploration, Resumes, Networking, Using LinkedIn Interviewing, The Job Search for International Students, Managing Offers and Negotiating.

\item \textbf{PSU Writing Center:} The Writing Center (GWC) provides free one-to-one peer consultations and interactive workshops for Penn State students of all disciplines and all levels of writing ability. The Writing Center provides several workshops each semester to help students learn about a variety of writing topics both in-person and online appointments for fifty (50) minutes per session.  The schedule varies each week. 

\item \textbf{Penn State Center for the Integration of Research, Teaching and Learning (CITRL):} The CIRTL Network offers a wide array of courses, workshops, and events primarily targeted to graduate students and postdocs but open to faculty as well on topics ranging from inclusive teaching to navigating mentoring and advising relationships. The mission of CIRTL is to enhance excellence in undergraduate education through the development of a national faculty committed to implementing and advancing effective teaching practices for diverse learners as part of successful and varied professional careers. Courses include Research Mentor Training, Laying the Foundations for a Successful Teaching Career, Advancing Learning Through Evidence-Based STEM Teaching, and Imaging Post-PhD Career Versatilty.

\item \textbf{PSU Learning Resource Network:} Penn State Faculty, Staff, Students and Volunteers can access content that covers topics including compliance, core business skills, project management, faculty development, technology, professional growth, and many more. Online courses in educational equity assist in creating a climate of diversity, equity, and inclusion; online courses in teaching excellence assist faculty in achieving teaching excellence by offering a variety of instructional tools and resources; online courses in research protections ensure that research is conducted in accordance with federal, state, local, and funding regulations.

\item \textbf{PSU Scholarship and Research Integrity Education Program} As research has become more complex, more collaborative, and more costly, issues of research ethics have become similarly complex, extensive, and important. The education of Penn State  students must prepare them to face these issues in their professional lives. The SARI (Scholarship and Research Integrity) program at Penn State is designed to offer graduate students comprehensive, multilevel training in the responsible conduct of research (RCR), in a way that is tailored to address the issues faced by students in individual programs.
All Penn State students must complete SARI requirements during their study.

\item \textbf{PSU Undergraduate Research and Fellowships Mentoring (URFM) Program} URFM believes that academic research should prioritize the mentorship of undergraduates. To this end, they have compiled resources on mentorship development and best practices. Additionally, they offer a recorded webinar on high-quality mentoring, which offers thoughtful reflection on undergraduate research mentoring which can be viewed by all project members.

\item \textbf{Community, State, and Federal Agency Collaborators:} 
The project will leverage established relationships with key local, state, and federal entities to ensure robust community engagement and structural alignment with national risk and preservation standards. 
\begin{itemize}
    \item \textbf{Department of the Interior / Jennifer Wellock:} Jennifer Wellock, co-author of the National Park Service flood adaptation guidelines, serves as the project's federal heritage recovery representative. She will provide expert oversight by reviewing completed building intervention typologies for consistency with federal practice and benchmarking the tabletop evaluation exercises. Her Department of the Interior connections will also facilitate nationwide dissemination of the project's toolkit to other heritage-at-risk communities.
    \item \textbf{Vermont State Historic Preservation Officer (SHPO):} The Vermont SHPO provides critical institutional context and historic preservation expertise for the Montpelier case study. They will participate as part of the expert panel evaluating pre- and post-disaster decision-making benchmarks during the annual APT conference tabletop exercises.
    \item \textbf{North Carolina State Historic Preservation Officer (SHPO):} The North Carolina SHPO provides matching regional institutional context and historic preservation data for the rural Appalachian case study in Marshall. They will join the expert panel alongside the Vermont SHPO to evaluate the baseline tabletop exercises and analyze the efficacy of the developed decision-support frameworks.
    \item \textbf{Montpelier Alive:} Operating as a vital downtown merchant and community network within the state capital, this organization will provide historic programmatic archives, email lists, and local meeting records to help empirically construct the district-scale collective action and exclusion patterns. Their community liaisons will also lead personalized door-to-door and mail outreach to ensure that historically marginalized property owners participate in the pre-disaster stakeholder workshops.
    \item \textbf{Foundation for a Resilient Montpelier:} Serving as a specialized local governance and recovery entity, the Foundation will provide localized disaster recovery records and institutional frameworks to evaluate municipal resilience capacity. They will assist the research team in translating spatial hydrodynamic risk outputs into actionable, municipal-scale policy proposals during the co-production phases.
    \item \textbf{Association for Preservation Technology Disaster Response Initiative (APT DRI):} As a primary practitioner partner, APT DRI—which PI Napolitano directs—serves as the formal platform for co-developing, field-testing, and hosting the Flood-Resilient Heritage Preservation Toolkit. The initiative will host a dedicated tabletop simulation exercise at its annual conference to empirically measure and score practitioner decision-making improvements using the toolkit.
    \item \textbf{U.S. Army Corps of Engineers (USACE):} To ensure the final toolkit effectively bridges the gap between historical preservation and municipal emergency management, the project incorporates flood risk communication frameworks from USACE. The research team will adapt resources from the USACE Flood Risk Communication Toolbox to establish trusted, multi-audience communication strategies tailored specifically to owners navigating the preservation-urgency paradox.
\end{itemize}


\end{itemize}
\end{document}